\documentclass[12pt]{article}
\usepackage[utf8]{inputenc}
\usepackage[T1]{fontenc}
\usepackage{hyperref}
\title{Seluwasan Language Resources for K-12 Curriculum Development}
\author{Compiled by Claude Sonnet 5.5 (claude-sonnet-5-5) for Andr\'as Kornai}
\date{}
\begin{document}
\maketitle

\section{Introduction}
Seluwasan (ISO 639-3 sws; Glottocode selu1242) is an Austronesian language of Yamdena in the Tanimbar Islands, Maluku, Indonesia. Search results give conflicting speaker figures: 2,840 (1980) on Wikipedia and Global Recordings, and 3,600 for the people group on Joshua Project. Joshua Project notes language shift under outside influence. Little material exists. No school textbooks, primers or dictionaries turned up in the searches.

\section{Primary Educational Resources}
No primers, textbooks or school programmes for Seluwasan were found in search results.

\section{Linguistic and Reference Materials}
\begin{itemize}
\item Wikipedia, ``Seluwasan language'' (overview, classification, dialect Makatian): \url{https://en.wikipedia.org/wiki/Seluwasan_language}
\item Wikipedia, ``Selaru languages'' (classification context): \url{https://en.wikipedia.org/wiki/Selaru_languages}
\end{itemize}

\section{Religious and Community Resources}
\begin{itemize}
\item ScriptureEarth page for sws, listing a Makatian Bible resource: \url{https://www.scriptureearth.org/00eng.php?iso=sws}
\item Joshua Project, Seluwasan people group (reports that Bible translation has started): \url{https://joshuaproject.net/people_groups/14836}
\item Joshua Project language page for sws: \url{https://joshuaproject.net/languages/sws}
\end{itemize}

\section{Digital and Audio Resources}
\begin{itemize}
\item Global Recordings Network, Seluwasan language page (audio recordings may be listed there; contents not verified here): \url{https://globalrecordings.net/en/language/sws}
\end{itemize}

\section{Curriculum Development Context}
The language is not known to be taught in schools. Any curriculum would need new materials built from fieldwork, since no pedagogical resources were found. The Makatian scripture material and the Global Recordings audio may serve as starting text and speech sources, after their content and licensing have been checked.

\end{document}